\section{定積分}
\begin{longtable}{@{}p{39mm}p{85mm}p{20mm}@{}}
\toprule
式に見えるもの & まず確かめること & 本文\\
\midrule
\endhead
対称な区間$[-a,a]$ & $f(-x)$を計算し,偶関数・奇関数に分けられるかを調べる. & \chapterref{章:定積分の対称性}\\
区間$[a,b]$と両端の対称性 & $x\mapsto a+b-x$で変形し,元の被積分関数と足すと簡単になるかを見る. & \chapterref{章:定積分の対称性}\\
$\sin x$と$\cos x$が対称 & $x\mapsto\frac\pi2-x$で分子と分母がどう変わるかを調べる. & \chapterref{章:定積分の対称性}\\
分母$1+\e^x$と偶関数の因子 & $\frac1{1+\e^x}+\frac1{1+\e^{-x}}=1$を使い,左右を合わせる. & \chapterref{章:定積分の対称性}\\
$k$と$n$を含む和 & 幅$\frac1n$と高さ$f(\frac kn)$を取り出す.区間も確認する. & \chapterref{章:積分と数列}\\
自然数$n$を含む積分 & 部分積分などで隣の添字と結ぶか,積分のまま上下を評価する. & \chapterref{章:積分と数列}\\
原始関数が見つからない & 値の表示・近似・大小比較・収束判定のどれが必要かを分ける. & \chapterref{章:近似と広義積分}\\
端点や内部で関数が発散 & 各特異点の片側を独立に切り取り,通常の積分を計算してから極限を取る. & \chapterref{章:近似と広義積分}\\
\bottomrule
\end{longtable}


